\centerline{\bf \Huge Моя хата с краю!}

\begin{flushright}
        \scriptsize{Первым трём решившим все задачи даю тыщу!}
\end{flushright}

\problem{1}
На сторонах выпуклого четырёхугольника как на диаметрах построены четыре круга. Докажите, что они покрывают весь четырёхугольник.

\problem{2}
На поле встретились 2025 ковбоев. Известно что расстояния между ковбоями попарно различны. Все ковбои одновременно выстрелили в ближайшего к ним (по расстоянию) ковбоя. Докажите, что кто-то останется в живых.

\problem{3}
На плоскости синим и красным цветом окрашено несколько точек так, что никакие три точки одного цвета не лежат на одной прямой (точек каждого цвета не меньше трёх). Докажите, что какие-то три точки одного цвета образуют треугольник, на трёх сторонах которого лежит не более двух точек другого цвета.


\problem{4}
В обоих пунктах вам нужно взять рандомную прямую (в пункте а) эта прямая должна быть не перпендикулярна ни одному из отрезков) и спроецируйте всё на эту прямую.

a) Можно ли на плоскости расположить 993 отрезка так, чтобы каждый отрезок обоими своими концами упирался строго внутрь каких-то двух других отрезков?

б) На плоскости дано конечное множество многоугольников (не обязательно выпуклых), каждые два из которых имеют общую точку. Докажите, что существует прямая, имеющая общие точки со всеми этими многоугольниками. 

\problem{5}
Квадрат разрезали на конечное число прямоугольников. Обязательно ли найдётся отрезок, соединяющий центры (точки пересечения диагоналей) двух прямоугольников, не имеющий общих точек ни с какими другими прямоугольниками, кроме этих двух?
